\section{A* Search}

A* is a best-first graph search algorithm that finds the least-cost path from an initial state to a goal state. \cite{hart1968formal} Unlike Breadth-First Search, which expands nodes uniformly by depth, A* prioritizes nodes using an evaluation function that combines the cost already incurred with an estimate of the remaining cost to the goal:

$$f(n) = g(n) + h(n)$$

where $g(n)$ is the exact cost of the path from the start to node $n$, and $h(n)$ is a heuristic estimate of the remaining cost from $n$ to the goal. A* maintains a priority queue sorted by $f(n)$, always expanding the lowest-cost node first.

The correctness of A* depends on the admissibility of the heuristic. A heuristic $h$ is \textit{admissible} if it never overestimates the true cost to the goal: $h(n) \leq h^*(n)$ for all $n$, where $h^*(n)$ is the true minimum cost. Under this condition, A* is guaranteed to find an optimal solution.

The practical performance of A* is determined largely by the quality of $h$. A tight heuristic directs the search along the optimal path with minimal node expansions, while a weak one causes the frontier to grow combinatorially. For complex combinatorial domains such as Sokoban, designing an admissible heuristic that accurately approximates true cost is difficult due to intricate interactions between boxes, walls, and goals. This motivates replacing the hand-crafted heuristic with a learned approximation via a deep neural network.

\section{DeepCubeA}

DeepCubeA, proposed by Agostinelli et al. \cite{agostinelli2019solving}, replaces the hand-crafted heuristic in A* with a deep neural network $h_\theta$ trained to estimate the cost-to-go from any state to the goal. The network is trained through a self-supervised procedure called \textit{Autodidactic Iteration}: beginning from the solved (goal) state, random sequences of actions of increasing length are applied to generate scrambled states. At each depth $d$, the network learns to predict $d$ as the cost-to-go from the resulting state. This backward induction scheme produces a heuristic that approximates true distance across the reachable state space without any hand-engineered domain knowledge.

At inference time, DeepCubeA uses Weighted A* (WA*), scaling the learned heuristic by a weight $w \geq 1$:

$$f(n) = g(n) + w \cdot h_\theta(n)$$

A weight $w > 1$ biases node expansion toward the goal, sacrificing the optimality guarantee in exchange for significantly faster convergence — a practical trade-off for large state spaces where optimal search is intractable. DeepCubeA demonstrated strong performance on challenging combinatorial puzzles including the Rubik's Cube and the 15-puzzle, establishing it as a general-purpose learned search framework. A subsequent work extending this approach, which uses Answer Set Programming to specify goals to deep neural networks, demonstrated the framework's applicability to Sokoban, successfully solving highly difficult instances involving more than 20 boxes. \cite{agostinelli2021admissible}
